\chapter{Multivariate random variables}

\begin{definition}[Bivariate random variable]
    Consider a probability space $(\Omega, \mathcal{F}, P)$ arising out of a 
    random experiment. A vector of functions $X = (X_1, X_2)$ that maps 
    $\Omega$ into $\mathbb{R}^2$ is said to be a bivariate random variable if for each $x_1, x_2 \in \mathbb{R}$, the set
    \begin{equation}
        \{ \omega \in \Omega \> | \> X_1(\omega) \leq x_1, X_2(\omega) \leq x_2 \} \in \mathcal{F}.
    \end{equation}
\end{definition}

\begin{example}
    An unbiased coin is tossed twice. Let $X_1$ be the number of heads obtained, and $X_2$ be the number of tails obtained. Then $X = (X_1, X_2)$ is a bivariate random variable 
    with sample space 
    \[
        \Omega = \{ (H,H), (H,T), (T,H), (T,T) \}
    \]
    where $H$ represents heads and $T$ represents tails. We take $\mathcal{F}$ to be the power set of $\Omega$.
    Notice that 
    \[
        \{ \omega \in \Omega \> | \> X_1(\omega) \leq x_1 \} = \begin{cases}
            \varnothing & x_1 < 0\\
            \{ (T,T) \} & 0 \leq x_1 < 1\\
            \{ (H,T), (T,H), (T,T) \} & 1 \leq x_1 < 2\\
            \Omega & x_1 \geq 2
        \end{cases}
    \]
    On the other hand, $X_2$ is defined similarly.
    \[
        \{ \omega \in \Omega \> | \> X_2(\omega) \leq x_2 \} = \begin{cases}
            \varnothing & x_2 < 0\\
            \{ (H,H) \} & 0 \leq x_2 < 1\\
            \{ (T,H), (H,T), (H,H) \} & 1 \leq x_2 < 2\\
            \Omega & x_2 \geq 2
        \end{cases}
    \]
\end{example}

\begin{theorem}
    Let $F_{X,Y}(x,y)$ be the distribution function of the jointly distributed random variables $X$, $Y$. Then
    \begin{enumerate}
        \item $F_{X,Y}(x, -\infty) = F_{X,Y}(-\infty, y) = 0$
        \item $F_{X,Y}(\infty, \infty) = 1$
        \item $F_{X,Y}(x + 0, y) = F_{X,Y}(x, y + 0) = F_{X,Y}(x, y)$
        \item For any $h > 0$ and $k > 0$,
        \[
            \Delta F_{X,Y}(x, y) = F_{X,Y}(x+h, y+k) - F_{X,Y}(x+h, y) - F_{X,Y}(x, y+k) + F_{X,Y}(x, y)
        \]
    \end{enumerate}
\end{theorem}
\begin{proof}
    \begin{enumerate}
        \item Let 
        \[A_n(x) = \{(X, Y): X \leq x, Y \leq -n\}\] 
        Notice that $A_{n+1}(x) \subseteq A_n(x)$ for all $n \in \mathbb{N}$, that is, $\{A_n\}$ is a monotone decreasing sequence of sets. So,
        \[
            \lim_{n \to \infty} A_n = \varnothing.
        \]
        Now,
        \begin{align*}
            F_{XY}(x, -\infty) &= \lim_{n \to \infty} \mathbb{P}[X \leq x, Y \leq -n]\\
            &= \lim_{n \to \infty} \mathbb{P}[A_n(x)]\\
            &= \mathbb{P}\left[ \lim_{n \to \infty} A_n(x) \right] && \text{(by continuity theorem of probability)}\\
            &= \mathbb{P}[\varnothing] = 0.
        \end{align*}
        Similarly, it can be shown that
        \[
            F_{XY}(-\infty, y) = \mathbb{P}[X \leq -\infty, Y \leq y] = 0.
        \]

        \item Let $A_n = \{(X, Y): -n \leq x \leq n, -n \leq y \leq n\}$.

        $A_n \subseteq A_{n+1}$, i.e., $\{A_n\}$ is a monotone expanding sequence of sets.
        \[
            \lim_{n \to \infty} A_n = \mathbb{R}^2.
        \]
        \begin{align*}
            F_{XY}(\infty, \infty) &= \lim_{n \to \infty} \mathbb{P}[X \leq n, Y \leq n] = \lim_{n \to \infty} \mathbb{P}(A_n) = \mathbb{P}\left[ \lim_{n \to \infty} A_n \right]\\
            &= \mathbb{P}[\mathbb{R}^2] && \text{(by continuity of probability)}
        \end{align*}
        Therefore, $F_{XY}(\infty, \infty) = 1$.

        \item[(iv)] $F_{XY}(x^+, y)$
        \[
            = \lim_{\varepsilon \to 0} F_{XY}(x + \varepsilon, y) = \lim_{n \to \infty} F_{XY}\left(x + \frac{1}{n}, y\right).
        \]
        Let $A_n = \{(X, Y): -\infty < X \leq x + \frac{1}{n}, -\infty < Y \leq y\}$.

        $A_n \supseteq A_{n+1}$, so
        \[
            \lim_{n \to \infty} A_n = \{(X, Y): -\infty < X \leq x, -\infty < Y \leq y\}.
        \]
        So,
        \begin{align*}
            \lim_{n \to \infty} F_{XY}\left(x + \frac{1}{n}, y\right) &= \lim_{n \to \infty} \mathbb{P}\left[-\infty < X \leq x + \frac{1}{n}, -\infty < Y \leq y\right]\\
            &= \lim_{n \to \infty} \mathbb{P}(A_n)\\
            &= \mathbb{P}\left[ \lim_{n \to \infty} A_n \right]\\
            &= \mathbb{P}[-\infty < X \leq x, -\infty < Y \leq y]\\
            &= F_{XY}(x, y).
        \end{align*}
        Similarly, it can be shown that
        \[
            F_{XY}(x, y^+) = F_{XY}(x, y).
        \]

        \item[(v)] $F_{XY}(x_1, y) \leq F_{XY}(x_2, y)$ and $F_{XY}(x, y_1) \leq F_{XY}(x, y_2)$.

        \item[(vi)] For every $x_1 > x_2$, $y_1 > y_2$,
        \[
            \Delta F_{XY}(x, y) = F_{XY}(x_1, y_1) - F_{XY}(x_1, y_2) - F_{XY}(x_2, y_1) + F_{XY}(x_2, y_2) \geq 0.
        \]
    \end{enumerate}
\end{proof}
\begin{proof}[Proof of (vi)]
    \begin{align*}
        \Delta F_{XY}(x, y) &= F_{XY}(x_1, y_1) - F_{XY}(x_1, y_2) - F_{XY}(x_2, y_1) + F_{XY}(x_2, y_2)\\
        &= \mathbb{P}[X \leq x_1, Y \leq y_1] - \mathbb{P}[X \leq x_1, Y \leq y_2]\\
        &\quad - \mathbb{P}[X \leq x_2, Y \leq y_1] + \mathbb{P}[X \leq x_2, Y \leq y_2]\\
        &= \mathbb{P}[X \leq x_1, y_2 < Y \leq y_1] - \mathbb{P}[X \leq x_2, y_2 < Y \leq y_1]\\
        &= \mathbb{P}[x_2 < X \leq x_1, y_2 < Y \leq y_1] \geq 0.
    \end{align*}
    Since $x_1 > x_2$, $y_1 > y_2$.
\end{proof}

\section{Conditional distributions}

\begin{definition}[Marginal probability density functions]
    If $X$ and $Y$ are two jointly continuous random variables with joint probability density function $F_{X,Y}(x,y)$, then the marginal probability density functions of $X$ is given by
    \begin{equation}
        F_X(x) = F_{X,Y}(x, \infty).
    \end{equation}
    Similarly, the marginal probability density function of $Y$ is 
    \begin{equation}
        F_Y(y) = F_{X,Y}(\infty, y).
    \end{equation}
\end{definition}

\begin{definition}[Conditional probability density function]
    Consider two jointly distributed discrete random variables $X$ and $Y$, then the conditional probability mass function of $Y$ given $X = x$ is defined as
    \begin{equation}
        F_{Y \mid X}(y \mid x) = \sum_{t \leq y} \mathbb{P}[Y = t \mid X = x] = \sum_{t \leq y} \frac{\mathbb{P}[X = x, Y = t]}{\mathbb{P}[X = x]}. \label{defn:con01}
    \end{equation}
\end{definition}

Recall the definition of conditional probabilities: For two sets 
$A$ and $B$, with $P(A) \neq 0$, the conditional probability of $B$ 
given that $A$ is true is defined as 

\begin{equation}
    \mathbb{P}(B \> | \> A) := \frac{\mathbb{P}(A \cap B)}{\mathbb{P}(A)}.
\end{equation}

From \cref{defn:con01}, we can said that the conditional probability mass function of $Y$ given $X = x$ is
\begin{align*}
    F_{Y \mid X}(y \mid x) &= \lim_{\varepsilon \to 0} \mathbb{P}[Y \leq y \mid x - \varepsilon \leq X \leq x + \varepsilon]\\
    &= \lim_{\varepsilon \to 0} \frac{\mathbb{P}[x - \varepsilon \leq X \leq x + \varepsilon, Y \leq y]}{\mathbb{P}[x - \varepsilon \leq X \leq x + \varepsilon]}.
\end{align*}
provided that the limit exists.

If there exists a non-negative function $f_{Y \mid X}(y \mid x)$ such that
\begin{equation}
    F_{Y \mid X}(y \mid x) = \int_{-\infty}^{y} f_{Y \mid X}(t \mid x) \, \mathrm{d}t,
\end{equation}
then $f_{Y \mid X}(y \mid x)$ is called the conditional probability density function of $Y$ given $X = x$.

\begin{remark}
    The conditional expectation of $Y$ given $X = x$ is called the 
    \textbf{regression} of $Y$ on $X$.
\end{remark}

\begin{example}
    Let $X$ and $Y$ be two jointly continuous random variable with joint density function
    \[
        f_{XY}(x, y) = \begin{cases}
            x^2 + \frac{1}{3}y, & -1 \leq x \leq 1,\> 0 \leq y \leq 1\\
            0 & \text{otherwise}
        \end{cases}.
    \]

    For $0 \leq y \leq 1$, find the conditional pdf of $X$ given $Y = y$.
\end{example}
\begin{solution}
    First we find the marginal distribution of $Y$, which we can obtain by integrating 
    along with $x$.
    \begin{align*}
        f_Y(y) = \int_{-1}^{1} f_{XY}(x,y)\> \mathrm{d}x &= 
        \int_{-1}^{1} \left( x^2 + \frac{1}{3}y \right) \> \mathrm{d}x\\
        &= \frac{1}{3}x^3 + \frac{1}{3}xy \bigg \vert^1_{-1}\\
        &= \frac{2}{3}(1+y).
    \end{align*}

    The conditional distribution of $X$ given $Y=y$ is 
    \[
        f_{X|Y=y}(x) = \frac{f_{XY}(x,y)}{f_Y(y)} 
        = \frac{x^2 + \mfrac{1}{3}y}{\mfrac{2}{3}(1+y)}
        = \frac{3x^2 + y}{2(1+y)}, \quad -1 \leq x \leq 1,\> 0 \leq y \leq 1
    \]
\end{solution}

\begin{example}
    Consider two random variables $X$ and $Y$ with joint probability density function
    \[
        f_{X,Y} (x,y) = \begin{cases}
            x^2 + \frac{1}{3}xy & 0 \leq x \leq 1, 0 \leq y \leq 2\\
            0 & \text{otherwise}
        \end{cases}
    \]
    Find
    \begin{enumerate}
        \item the marginal probability density functions of $X$,
        \item the marginal probability density functions of $Y$,
        \item the conditional probability density function of $Y < \frac{1}{2}$ given $X<\frac{1}{2}$.
    \end{enumerate}
\end{example}
\begin{solution}
    \begin{enumerate}
        \item The marginal probability density function of $X$ is 
        \begin{align*}
            f_X(x) &= \int^2_0 f_{X,Y}(x,y) \, \mathrm{d}y = 
            \int^2_0 \left( x^2 + \frac{1}{3}xy \right) \, \mathrm{d}y = \left[ x^2 y + \frac{1}{6}xy^2 \right]^2_0 = 2x^2 + \frac{2x}{3}, \quad 0 \leq x \leq 1.
        \end{align*}
        and \[\int^1_0 f_X(x) \> \mathrm{d}x = \int_{0}^{1} \left(2x^2 + \frac{2x}{3}\right) \mathrm{d}x = 1. \]
        So this is a valid probability density function.
        \item The conditional probability density function of $Y < \frac{1}{2}$ given $X<\frac{1}{2}$ is 
    \end{enumerate}
\end{solution}

\begin{example}
    Let $(X, Y)$ be a bivariate discrete random variable with joint PMF
    \[
        f_{XY}(x,y) = \begin{cases}
            \dfrac{1}{\binom{m+1}{2}} & y = 1, 2, \ldots, \alpha \text{ and } x = y, y+1, \ldots, m\\[6pt]
            0 & \text{otherwise}
        \end{cases}
    \]
    where $m$ is a positive integer $\geq 1$. Find $\mathbb{E}(X)$, $\mathbb{E}(X \mid Y)$, and $\mathbb{E}[\mathbb{E}(X \mid Y)]$.
\end{example}
\begin{solution}
    The marginal PMF of $X$ is
    \[
        f_X(x) = \sum_{y=1}^{x} \frac{1}{\binom{m+1}{2}} = \frac{x}{\binom{m+1}{2}}.
    \]
    Therefore,
    \[
        \mathbb{E}(X) = \sum_{x=1}^{m} x \cdot \frac{x}{\binom{m+1}{2}} = \frac{1}{\binom{m+1}{2}} \cdot \frac{m(m+1)(2m+1)}{6} = \frac{2m+1}{3}.
    \]
    The marginal PMF of $Y$ is
    \[
        f_Y(y) = \sum_{x=y}^{m} f_{XY}(x,y) = \frac{m-y+1}{\binom{m+1}{2}}.
    \]
    The conditional PMF of $X$ given $Y = y$ is
    \[
        f_{X \mid Y}(x \mid y) = \frac{f_{XY}(x,y)}{f_Y(y)} = \frac{\frac{1}{\binom{m+1}{2}}}{\frac{m-y+1}{\binom{m+1}{2}}} = \frac{1}{m-y+1}, \quad x = y, y+1, \ldots, m.
    \]
    Thus,
    \[
        \mathbb{E}[X \mid Y = y] = \sum_{x=y}^{m} x \cdot \frac{1}{m-y+1} = \frac{1}{m-y+1} \cdot (m-y+1) \cdot \frac{m+y}{2} = \frac{m+y}{2}.
    \]
    Now,
    \begin{align*}
        \mathbb{E}[\mathbb{E}(X \mid Y)] &= \mathbb{E}\left[ \frac{m+Y}{2} \right]\\
        &= \frac{m}{2} + \frac{1}{2} \mathbb{E}(Y)\\
        &= \frac{m}{2} + \frac{1}{2} \left\{ \sum_{y=1}^{m} y \cdot \frac{m-y+1}{\binom{m+1}{2}} \right\}\\
        &= \frac{m}{2} + \frac{1}{2} \cdot \frac{1}{\binom{m+1}{2}} \left[ \frac{m \cdot m(m+1)}{2} + \frac{m(m+1)}{2} - \frac{m(m+1)(2m+1)}{6} \right]\\
        &= \frac{m}{2} + \left[ \frac{m}{2} + \frac{1}{2} - \frac{2m+1}{6} \right] = \frac{2m+1}{3}.
    \end{align*}
    This verifies the law of iterated expectations: $\mathbb{E}[\mathbb{E}(X \mid Y)] = \mathbb{E}(X)$.
\end{solution}

\begin{example}
    Let $X$ and $Y$ be two jointly distributed continuous random variables with joint PDF
    \[
        f_{X,Y}(x,y) = \frac{1}{2\pi\sqrt{1-\rho^2}} \exp\left\{ -\frac{1}{2(1-\rho^2)} \left[ (x - \mu_X)^2 - 2\rho(x-\mu_X)(y-\mu_Y) + (y-\mu_Y)^2 \right] \right\}, \quad y \in \mathbb{R}.
    \]
    \begin{enumerate}
        \item Find the marginal PDF of $X$.
        \item Find the conditional PDF of $Y$ given $X = x$.
    \end{enumerate}
\end{example}
\begin{solution}
    \begin{enumerate}
        \item The marginal PDF of $X$ is
        \begin{align*}
            f_X(x) &= \int_{-\infty}^{\infty} f_{X,Y}(x,y) \, \mathrm{d}y \\
            &= \int_{-\infty}^{\infty} \frac{1}{2\pi\sqrt{1-\rho^2}} \exp\left\{ -\frac{1}{2(1-\rho^2)} \left[ (x-\mu_X)^2 - 2\rho(x-\mu_X)(y-\mu_Y) + (y-\mu_Y)^2 \right] \right\} \mathrm{d}y \\
            &= \frac{e^{-\frac{(x-\mu_X)^2}{2}}}{2\pi\sqrt{1-\rho^2}} \int_{-\infty}^{\infty} \exp\left\{ -\frac{\left[ (y-\mu_Y) - \rho(x-\mu_X) \right]^2}{2(1-\rho^2)} \right\} \mathrm{d}y.
        \end{align*}
        Using the fact that $\displaystyle\int_{-\infty}^{\infty} \exp\left\{ -\frac{(z-c)^2}{2\sigma^2} \right\} \mathrm{d}z = \sigma\sqrt{2\pi}$, we have
        \[
            f_X(x) = \frac{1}{\sqrt{2\pi}} \exp\left\{ -\frac{(x-\mu_X)^2}{2} \right\}, \quad x \in \mathbb{R}.
        \]
        Hence, $X \sim N(\mu_X, 1)$.

        \item The conditional PDF of $Y$ given $X = x$ is
        \begin{align*}
            f_{Y \mid X}(y \mid x) &= \frac{f_{X,Y}(x,y)}{f_X(x)} \\
            &= \frac{1}{\sqrt{2\pi(1-\rho^2)}} \exp\left\{ -\frac{\left[ (y - \mu_Y) - \rho(x - \mu_X) \right]^2}{2(1-\rho^2)} \right\}.
        \end{align*}
        Therefore, $Y \mid X = x \sim N\left( \mu_Y + \rho(x - \mu_X), 1 - \rho^2 \right)$.

        The conditional expectation is
        \[
            \mathbb{E}[Y \mid X = x] = \mu_Y + \rho(x - \mu_X).
        \]
        And the unconditional expectation is
        \[
            \mathbb{E}(Y) = \int_{-\infty}^{\infty} y \cdot f_Y(y) \, \mathrm{d}y = \mu_Y.
        \]
    \end{enumerate}
\end{solution}

\begin{example}
    An exam consists of a problem section and a short-answer. Let $X$ denote the amount of time (in hours) a student spends on the problem section, and let $Y$ denote the amount of time (in hours) the same 
    student spends on the short-answer section. Suppose that the joint PDF of $X$ and $Y$ is given by
    \[
        f_{X, Y}(x, y) = \begin{cases}
            \displaystyle \frac{432}{5}x^3y & 0 \leq x \leq 1, \frac{x}{3} \leq y \leq \frac{x}{2}.\\
            0 & \text{otherwise}
        \end{cases}
    \]
    \begin{enumerate}
        \item Find the probability that a student spends at most 0.25 hours on the short-answer section.
        \item If the student spends exactly 0.25 hour on the short-answer section, what is the 
        probability that at most 0.6 hour was spent on the problem section?
        \item Find the average difference of the amount of time spent on the problem section and 
        short-answer section of the exam.
    \end{enumerate}
\end{example}
\begin{solution}
    Given the joint PDF
    \[
        X,Y \sim f_{X,Y}(x,y) = \begin{cases}
            \displaystyle \frac{432}{5} x^3 y & 0 \leq x \leq 1, \frac{x}{3} \leq y \leq \frac{x}{2}\\
            0 & \text{otherwise}
        \end{cases}
    \]
    \begin{enumerate}
        \item The probability that a student spends at most 0.25 hours on the short-answer section is
        \begin{align*}
            \mathbb{P}[Y > 0.25] &= 1 - \mathbb{P}[Y \leq 0.25]
            = 1 - \int^{1/4}_0 \int^{3y}_{2y} \frac{432}{5} x^3 y \, \mathrm{d}x \, \mathrm{d}y\\
            &= 1 - \int^{1/4}_0 \frac{432}{20} x^4y \Big|^{x=3y}_{x=2y} \, \mathrm{d}y\\
            &= 1 - \int^{1/4}_0 \left( \frac{432}{5}y \cdot \frac{65}{4}y^4 \right) \, \mathrm{d}y\\
            &= 1 - \int^{1/4}_0 1404y^5 \, \mathrm{d}y\\
            &= 1 - \frac{1404}{6}y^6 \Big|^{y=1/4}_{y=0}\\
            &= \fbox{$0.9429$}.
        \end{align*}

        \item The marginal probability of $X$ where $Y = 0.25$ is 
        \begin{align*}
            \int^{3/4}_{1/2} f_{X \mid Y = 0.25} (x, 0.25) \, \mathrm{d}x 
            &= \int^{3/4}_{1/2} \frac{432}{5} \cdot 0.25x^3 \, \mathrm{d}x\\
            &= \frac{432}{5} \cdot 0.25 \cdot \frac{x^4}{4} \Big|^{x=3/4}_{x=1/2}\\
            &= 1.3711
        \end{align*}

        Given $Y =0.25$, the probability that at most 0.6 hour was spent on the problem section is
        \begin{align*}
            \mathbb{P}[X \leq 0.6 \mid Y = 0.25] &= \int^{0.6}_{0.5} f_{X \mid Y = 0.25} (x, 0.25) \, \mathrm{d}x
            = \frac{\displaystyle \int^{0.6}_{0.5} \frac{432}{5} \cdot 0.25x^3 \, \mathrm{d}x}{\displaystyle \int^{3/4}_{1/2} f_{X \mid Y = 0.25}(x, 0.25) \mathrm{d}x}\\
            &= \frac{108}{5 (1.3711)} \cdot \left( \frac{x^4}{4} \right)\Big|^{x=0.6}_{x=0.5}\\[0.7em]
            &= \fbox{$0.2643$}.
        \end{align*}

        \item Compute the expected value of $X - Y$:
        \begin{align*}
            \mathbb{E}[X - Y] &= \mathbb{E}[X] - \mathbb{E}[Y]\\
            &= \int^1_0 \int^{x/2}_{x/3} x \cdot \frac{432}{5} x^3 y \, \mathrm{d}y \, \mathrm{d}x - 
            \left[ \int_{0}^{1/3} \int^{3y}_{2y} y \left( \frac{432}{5} x^3 y \right) \, \mathrm{d}x \, \mathrm{d}y 
            + \int_{1/3}^{1/2} \int_{2y}^{1} y \left( \frac{432}{5} x^3 y \right) \, \mathrm{d}x \, \mathrm{d}y \right]\\
            &= \int^1_0 \frac{432}{5} \cdot \frac{1}{2}x^4 y^2 \Big|^{y=x/2}_{y=x/3} \, \mathrm{d}x - \left[
                1404 \int^{1/3}_{0} y^6 \, \mathrm{d}y + \frac{432}{5} \int_{1/3}^{1/2} y^2 \left( \frac{1}{4} - \frac{(2y)^4}{4} \,  \right)\mathrm{d}y
            \right]\\
            &= \int_{0}^{1} 6x^6 \, \mathrm{d}x - \left[
                1404 \cdot \frac{y^7}{7} \Big|^{y=1/3}_{y=0} + \left( \frac{36}{5}y^3 - \frac{1728}{35}y^7 \right) \Big|^{y=1/2}_{y=1/3}
            \right]\\
            &= \frac{6}{7} - \left(\frac{4}{35} + \frac{26}{105}\right)\\ 
            &= \fbox{$\displaystyle \frac{52}{105}$} \quad \text{ hour}.
        \end{align*}
        Thus, the average difference of the amount of time spent on the problem section and short-answer section of the exam is
        $52 / 105 \approx 0.495$ hours.
    \end{enumerate}
\end{solution}

\section{Multivariate Normal Distribution}

An absolutely continuous random vector $(X_1, X_2)$ is said to follow a bivariate normal distribution with parameters $(\mu_1, \mu_2, \sigma_1, \sigma_2, \rho)$ if the joint PDF of $(X_1, X_2)$ is
\[
    f(x_1, x_2) = \frac{1}{\sigma_1 \sigma_2 \cdot 2\pi \sqrt{1 - \rho^2}} \, e^{-\frac{1}{2} \varphi(x_1, x_2)}, \quad x_1, x_2 \in \mathbb{R}, \; \sigma_1, \sigma_2 \in \mathbb{R}, \; \mu_1, \mu_2 \in \mathbb{R}, \; |\rho| < 1
\]
where
\[
    \varphi(x_1, x_2) = \frac{1}{1 - \rho^2} \left[ \left( \frac{x_1 - \mu_1}{\sigma_1} \right)^2 - 2\rho \left( \frac{x_1 - \mu_1}{\sigma_1} \right) \left( \frac{x_2 - \mu_2}{\sigma_2} \right) + \left( \frac{x_2 - \mu_2}{\sigma_2} \right)^2 \right].
\]
We write $(X_1, X_2) \sim N(\mu_1, \mu_2, \sigma_1, \sigma_2, \rho)$.

The marginal PDF of $X_1$ is
\begin{align*}
    f_{X_1}(x_1) &= \int_{-\infty}^{\infty} f(x_1, x_2) \, \mathrm{d}x_2 = \int_{-\infty}^{\infty} \frac{1}{\sigma_1 \sigma_2 2\pi \sqrt{1 - \rho^2}} \, e^{-\frac{1}{2} \varphi(x_1, x_2)} \, \mathrm{d}x_2
\end{align*}
where
\[
    \varphi(x_1, x_2) = \frac{1}{1 - \rho^2} \left[ \left( \frac{x_2 - \mu_2}{\sigma_2} \right) - \rho \left( \frac{x_1 - \mu_1}{\sigma_1} \right) \right]^2 + \frac{1}{\sigma_1^2} \left( x_1 - \mu_1 \right)^2.
\]
Therefore,
\begin{align*}
    f_{X_1}(x_1) &= \int_{-\infty}^{\infty} \frac{1}{2\pi \sigma_1 \sigma_2 \sqrt{1 - \rho^2}} \, e^{-\frac{1}{2(1-\rho^2)} \left[ \left( \frac{x_2 - \mu_2}{\sigma_2} \right) - \rho \left( \frac{x_1 - \mu_1}{\sigma_1} \right) \right]^2 - \frac{1}{2} \left( \frac{x_1 - \mu_1}{\sigma_1} \right)^2} \, \mathrm{d}x_2 \\
    &= \frac{1}{\sigma_1 \sqrt{2\pi}} \, e^{-\frac{1}{2} \left( \frac{x_1 - \mu_1}{\sigma_1} \right)^2}, \quad -\infty < x_1 < \infty
\end{align*}

since the integrand is in the form of a normal PDF. Similarly, it can be shown that
\[
    f_{X_2}(x_2) = \frac{1}{\sigma_2 \sqrt{2\pi}} \, e^{-\frac{1}{2} \left( \frac{x_2 - \mu_2}{\sigma_2} \right)^2}, \quad -\infty < x_2 < \infty.
\]

\begin{theorem}
    If the regression of $Y$ on $X$ is linear and the variance of $Y$ is algebraically independent of $X$, then
    \[
        V(Y \mid X) = \sigma_Y^2 (1 - \rho^2).
    \]
\end{theorem}
\begin{proof}
    \begin{align*}
        V(Y \mid X) &= \mathbb{E}\left[ (Y - \eta_X)^2 \mid X \right]\\
        &= \mathbb{E}\left[ (Y - \alpha - \beta X)^2 \mid X \right] && \text{Regression is linear}\\
        &= \mathbb{E}\left[ (Y - \mu_Y) - \beta (X - \mu_X) \right]^2 && \text{Since it is independent of } X\\
        &= \mathbb{E}\left[ \left\{ (Y - \mu_Y) - \rho \frac{\sigma_Y}{\sigma_X} (X - \mu_X) \right\}^2 \right]\\
        &= \mathbb{E}(Y - \mu_Y)^2 - 2\rho \frac{\sigma_Y}{\sigma_X} \mathbb{E}(Y - \mu_Y)(X - \mu_X) + \rho^2 \frac{\sigma_Y^2}{\sigma_X^2}\, \mathbb{E}(X - \mu_X)^2\\
        &= \sigma_Y^2 - 2\rho \frac{\sigma_Y}{\sigma_X} \times \rho \sigma_X \sigma_Y + \rho^2 \frac{\sigma_Y^2}{\sigma_X^2} \times \sigma_X^2\\
        &= \sigma_Y^2 - 2\rho^2 \sigma_Y^2 + \rho^2 \sigma_Y^2\\
        &= \sigma_Y^2 (1 - \rho^2).
    \end{align*}
\end{proof}

\subsection{Correlation index and Correlation ratio}

When a regression curve is fitted to the conditional mean, we may denote by $\hat{Y}_X$ the estimated value of $\eta_X$ and by $\hat{\varepsilon}$ the residual $Y - \hat{Y}_X$. The variance of $\hat{\varepsilon}$ may be taken as an index of the usefulness of the fitted curve. Note in case, where a linear regression equation is fitted on $\eta_X$.

Then
\begin{align*}
    V(\hat{\varepsilon}) &= V(Y - \hat{Y}_X) = V(Y - \alpha - \beta X)\\
    &= V\left[ (Y - \mu_Y) - \rho \frac{\sigma_Y}{\sigma_X} (X - \mu_X) \right]\\
    &= V(Y - \mu_Y) + \rho^2 \frac{\sigma_Y^2}{\sigma_X^2} V(X - \mu_X) - 2\rho \frac{\sigma_Y}{\sigma_X} \text{Cov}\left[ (X - \mu_X), (Y - \mu_Y) \right]\\
    &= \sigma_Y^2 + \rho^2 \frac{\sigma_Y^2}{\sigma_X^2} \cdot \sigma_X^2 - 2\rho \frac{\sigma_Y}{\sigma_X} \cdot \rho \sigma_X \sigma_Y\\
    &= \sigma_Y^2 (1 - \rho^2).
\end{align*}

Now, $\text{Cov}(Y, \hat{\varepsilon}) = 0$ from the normal equation. Therefore,
\begin{align*}
    \text{Cov}(Y, Y - \hat{Y}_X) &= 0\\
    \text{Cov}(Y, \hat{Y}_X) &= V(Y).
\end{align*}
Hence,
\[
    V(\hat{\varepsilon}) = V(Y) - V(\hat{Y}_X).
\]

Now,
\[
    \rho^2_{Y\hat{Y}} = 1 - \frac{V(\hat{\varepsilon})}{V(Y)}, \quad \text{or} \quad \rho^2 = 1 - \frac{V(\hat{\varepsilon})}{V(Y)} = \frac{V(\hat{Y}_X)}{V(Y)} = \frac{V(\hat{Y})}{V(Y)}.
\]

\section{Regression}

Consider two random variables $X$ and $Y$, where $Y$ is the study of variance and $X$ is 
the auxilary variable. Our problem is to predict $Y$, when $X$ is known. 
Let $\eta_X = \mathbb{E}[Y \mid X = x]$ predict at random for the prediction of the form $\phi(X)$ to minimize $\mathbb{E}[(Y - \phi(X))^2]$. Then this is minimum when $\phi(X) = \eta_X$, the conditional mean.

Let $g(X) = \eta_X$, where $g(x)$ is any prediction of $Y$ based on $X$. Then
\begin{align*}
    \mathbb{E}[(Y - g(X))^2] &= \mathbb{E}\left\{ \mathbb{E}[(Y - g(X))^2 \mid X] \right\}\\
    &= \mathbb{E}\left\{ \mathbb{E}[(Y - \eta_X + \eta_X - g(X))^2 \mid X] \right\}\\
    &= \mathbb{E}\left\{ \mathbb{E}[(Y - \eta_X)^2 \mid X] + (\eta_X - g(X))^2 + 2(Y - \eta_X)(\eta_X - g(X)) \right\}\\
    &= \mathbb{E}\left\{ \mathbb{E}[(Y - \eta_X)^2 \mid X] \right\} + \mathbb{E}[(\eta_X - g(X))^2].
\end{align*}

$\mathbb{E}[(Y - g(X))^2]$ is minimum when $g(X) = \eta_X$.

Now,
\begin{align*}
    \mathbb{E}[(Y - \eta_X) \cdot (g(X) - \eta_X)] &= \mathbb{E}[(g(X) - \eta_X) \cdot \mathbb{E}[(Y - \eta_X) \mid X]]\\
    &= \mathbb{E}[(g(X) - \eta_X) \times 0] = 0.
\end{align*}

\subsection{Correlation coefficient between $Y$ and $\eta_X$}

\begin{align*}
    \text{Cov}(Y, \eta_X) &= \mathbb{E}[(Y - \mathbb{E}(Y))(\eta_X - \mathbb{E}(\eta_X))]\\
    &+ \mathbb{E}[(Y - \mathbb{E}(Y))(\eta_X - \mathbb{E}(Y))]\\
    &= \mathbb{E}\left[ (\eta_X - \mathbb{E}(Y)) \mathbb{E}[(Y - \mathbb{E}(Y)) \mid X] \right]\\
    &= \mathbb{E}[\eta_X - \mathbb{E}(Y)]^2\\
    &= V(\eta_X).
\end{align*}

Therefore,
\[
    \rho_{Y, \eta_X} = \frac{V(\eta_X)}{\sqrt{V(Y) \cdot V(\eta_X)}} = \sqrt{\frac{V(\eta_X)}{V(Y)}}.
\]

Hence, $\rho_{Y, \eta_X}^2 \geq \rho_{Y, g(X)}^2$ for any other prediction $g(X)$.

\begin{align*}
    \text{Cov}(Y, g(X)) &= \mathbb{E}[(Y - \mathbb{E}(Y))(g(X) - \mathbb{E}(g(X)))]\\
    &= \mathbb{E}\left[ (g(X) - \mathbb{E}(g(X))) \cdot \mathbb{E}[(Y - \mathbb{E}(Y)) \mid X] \right]\\
    &= \text{Cov}(\eta_X, g(X)).
\end{align*}

Therefore,
\[
    \rho_{Y, g(X)}^2 = \frac{\text{Cov}^2(Y, g(X))}{V(Y) \cdot V(g(X))} = \frac{\text{Cov}^2(\eta_X, g(X))}{V(Y) \cdot V(g(X))}.
\]

Hence, $\rho_{Y, g(X)} \leq \rho_{Y, \eta_X}$. This holds when $\eta_X$ and $g(X)$ have a linear relationship.

\begin{remark}
    The square of the maximum correlation is called the \textbf{correlation ratio} of $Y$ on $X$ and is usually denoted by
    \[
        \eta_{YX}^2 = \rho_{Y, \eta_X}^2, \quad \text{or} \quad \eta_{YX} \leq 1.
    \]
    It can be shown that
    \[
        V(Y) = V(\mathbb{E}[Y \mid X]) + \mathbb{E}[V(Y \mid X)] \geq V(\eta_X),
    \]
    since $V(Y \mid X) \geq 0$. Therefore,
    \[
        \frac{V(\eta_X)}{V(Y)} \leq 1 \quad \Rightarrow \quad \eta_{YX} \leq 1.
    \]
\end{remark}

\subsection{Determination of regression equation or regression curve}

Suppose the regression equation of $Y$ on $X$ is linear. Let $\eta_X = a + bX$. Then
\[
    \mathbb{E}(\eta_X) = a + b\mathbb{E}(X),
\]
so $\mu_Y = a + b\mu_X$.

Now,
\begin{align*}
    \mathbb{E}[X \eta_X] &= a\mu_X + b\mathbb{E}(X^2),\\
    \mathbb{E}(XY) &= a\mu_X + b(\sigma_X^2 + \mu_X^2).
\end{align*}

Therefore,
\begin{align*}
    \text{Cov}(X, Y) &= a\mu_X + b\sigma_X^2 + b\mu_X^2 - \mu_X(a + b\mu_X) = b\sigma_X^2.
\end{align*}

Hence,
\[
    b = \frac{\text{Cov}(X, Y)}{\sigma_X^2},
\]
and the regression equation of $Y$ on $X$ is
\[
    \eta_X = \mu_Y + \frac{\text{Cov}(X, Y)}{\sigma_X^2}(X - \mu_X).
\]

